\section{Limitations}
\label{sec:limitations}

This study has several limitations. Firstly, all results come from computer simulations. We did not make or measure any physical coatings. Therefore, the results show what the rescue method can achieve under our model. But we did not show how it would perform in an actual coating process.

Further, our model includes only thickness errors. Real coatings may have other problems, such as errors in refractive index, material absorption, and variations in layer thickness across the mirror surface. We did not include these effects. We also used the same refractive index for each material at 532~nm and 1064~nm. Real silica and tantala have slightly different refractive indices at these wavelengths. This could change both the original designs and the results of the rescue.

We also simplified the pass test by checking only the transmission at 1064~nm and 532~nm. The real coating specification also considers the electric field at the coating surface and how uniform the coating is across the mirror. We did not consider these requirements in our model.

Our assumptions about measurement errors may also differ from a real coating machine. We tested measurement errors from 0\% to 2\%, but these values were not measured from an actual monitoring system. We assumed that measurement errors were independent from layer to layer. In practice, they may be correlated or build up during a coating run. We also tested two separate deposition-error models: fully independent errors and one shared error across all layers. Real coating processes may contain both types at the same time. In addition, we did not include the fixed part of the deposition error, which can be about 1~nm per layer.

The rescue itself was tested in a limited way. We applied it once before the final few layers were deposited. We did not test applying it after every layer. We also did not test a method that estimates a shared deposition error from the measured layers and corrects the deposition rate.

Most of our results use four redesigned layers. We chose this number based on preliminary tests on one design, rather than a comparison across all 32 designs. Later tests showed that redesigning six or eight layers could improve the pass rate further. However, these tests covered only two error conditions: 2\% independent error and 1\% shared error. We therefore cannot yet say how many layers should be redesigned under different conditions. Similarly, our selective rescue method used the same imperfect measurements to decide whether a mirror needed rescue. A better decision rule that accounts for measurement uncertainty might perform better.

We also did not test whether all redesigned layers would be practical to make. We limited their thicknesses to between 0.5 and 1.5 times their planned values, and the redesigned layers received deposition errors just like the other layers. However, we did not examine how often the optimizer reached these limits or how difficult the new thicknesses would be to deposit in practice.

The results are based on one type of mirror, with two coating materials. We also tested 32 designs generated using the same optimization method and cost function. Three designs came from an earlier run whose original script was not saved. Removing these three designs did not change the main result, but designs generated using a different method may behave differently. We also did not compare the rescue method with designs that were made more robust against thickness errors from the beginning.

Each simulation point used 200 mirrors. This gives some uncertainty in the pass rates, and the uncertainty in an individual tolerable-error value is about 0.1--0.3 percentage points. The medians across all 32 designs are more stable, and a second set of simulated mirrors gave similar results. Still, more simulations would give more precise estimates.

Finally, our fast optimization method does not always find the best possible rescue. In a comparison with a slower, global optimizer, the slower method rescued some mirrors that the fast method missed. However, when we repeated this comparison on all 32 designs at 1\% shared error, the pass rates changed by no more than 3.5 percentage points. The noise proxy used during optimization is also based on physical layer thickness rather than optical thickness, as used in the LIGO specification. Its weight depends on the material values in the software library. We used the full thermal-noise calculation for the reported noise results, but the choice of optimization weights may still affect which designs and rescued coatings are selected.

These limitations mean that the results are evidence that the rescue method can work under certain simulated conditions, not proof that it is ready for use in a real coating process. Further simulations with more realistic error models, followed by experimental tests, are needed to establish its practical value.
